\documentclass[11pt]{article}

% ----------------------------------------------------------------------
%  Attracting the Elite
%  Coordinated Sorting of a Heterogeneous Tourist Population
%  under an Asymmetric Peer-Congestion Externality
%
%  Synthesis of two independently developed derivations
%  (Claude model.md + Codex model-codex.md), cross-checked.
% ----------------------------------------------------------------------

\usepackage[utf8]{inputenc}
\usepackage[T1]{fontenc}
\usepackage{lmodern}
\usepackage[margin=1in]{geometry}
\usepackage{amsmath,amssymb,amsthm}
\usepackage{mathtools}
\usepackage{bm}
\usepackage{booktabs}
\usepackage{enumitem}
\usepackage[hidelinks]{hyperref}

% --- theorem environments ---
\theoremstyle{plain}
\newtheorem{proposition}{Proposition}
\newtheorem{lemma}{Lemma}
\newtheorem{corollary}{Corollary}
\theoremstyle{definition}
\newtheorem{assumption}{Assumption}
\newtheorem{definition}{Definition}
\theoremstyle{remark}
\newtheorem{remark}{Remark}

% --- notation macros ---
\newcommand{\R}{\mathbb{R}}
\newcommand{\mH}{m_{H}}
\newcommand{\mL}{m_{L}}
\newcommand{\vH}{v_{H}}
\newcommand{\vL}{v_{L}}
\newcommand{\aHH}{\alpha_{HH}}
\newcommand{\aHL}{\alpha_{HL}}
\newcommand{\aLH}{\alpha_{LH}}
\newcommand{\aLL}{\alpha_{LL}}
\newcommand{\WP}{W^{P}}
\newcommand{\WS}{W^{S}}
\newcommand{\TP}{T^{P}}
\newcommand{\PiP}{\Pi^{P}}
\newcommand{\PiS}{\Pi^{S}}
\newcommand{\wHP}{w_{H}^{P}}
\newcommand{\wLP}{w_{L}^{P}}
\DeclarePairedDelimiter{\abs}{\lvert}{\rvert}
\newcommand{\pos}[1]{\left(#1\right)_{+}}

\title{\textbf{Attracting the Elite}\\[0.4em]
\large Coordinated Sorting of a Heterogeneous Tourist Population\\
under an Asymmetric Peer--Congestion Externality}
\author{Attracting-the-Elite project}
\date{June 2026}

\begin{document}
\maketitle

\begin{abstract}
\noindent
We study $J$ \emph{intrinsically identical} resorts facing a heterogeneous tourist
population of two types: rich/high--value $(H)$ and poor/low--value $(L)$. Payoffs are
composition--dependent: a visitor's utility falls in the mass of co--visitors of each type,
and the central structural assumption is that the externality is \emph{asymmetric}, the rich
strongly averse to crowding by the poor while the poor are nearly indifferent to the rich,
$\aHL \gg \aLH$. We ask whether resorts that are otherwise identical can raise their
\emph{joint} profit by coordinating on \emph{differentiated} prices that split the population
by type, relative to a symmetric baseline in which all resorts post one price. The answer is
yes, under a precise condition. We give an exact profit decomposition,
$\PiS - \PiP = \Delta W + \TP - R$, separating three forces: cross--type congestion relief,
own--type concentration cost, and the trade of the rent a common price must leave (the
pooling rent $\TP$) for the screening rent $R$ required to separate types. The asymmetry
plays a \emph{dual} role: $\aHL$ both creates real surplus and relaxes the high type's
incentive constraint, which is the screening technology. We also show a countervailing force
that informal accounts miss: a \emph{low} $\aLH$ makes the poor harder to keep out of premium
resorts. At the optimal split, sorting raises total surplus iff
$\aHL + \aLH > 2\sqrt{\aHH\,\aLL}$, and the same inequality governs the local instability of
the pooled equilibrium, which links the profit result to a wrong--equilibrium trap.
\end{abstract}

\tableofcontents
\bigskip

% ======================================================================
\section{Introduction}
% ======================================================================

Consider five ski resorts on comparable mountains, with comparable lifts, snow, and service.
Nothing physical distinguishes them. Yet we routinely observe such clusters \emph{sorting}
into a premium destination that serves a wealthy clientele at high prices and budget
destinations that serve a price--sensitive mass market. Where does the differentiation come
from when the product is identical?

This note isolates one mechanism: an \emph{asymmetric peer externality}. Wealthy tourists
dislike sharing a mountain with a large low--budget crowd, whereas low--budget tourists are
largely indifferent to the presence of the wealthy, and may even be drawn to it. Under this
asymmetry, a group of identical resorts can raise total profit by coordinating on a
\emph{menu} of prices that sorts the two types across locations, even though no resort has any
intrinsic advantage. The differentiation is created endogenously by the composition of who
shows up.

We keep three forces distinct throughout.
\begin{enumerate}[label=(\roman*),leftmargin=2.4em]
\item \textbf{Surplus creation.} Sorting physically separates the rich from the poor, deleting
the cross--congestion cost that the rich bear under mixing. Because the poor barely mind being
moved away from the rich, little or no offsetting loss arises, so the relief is a genuine
surplus increase.
\item \textbf{Screening.} For sorting to be self--enforcing, the rich must \emph{voluntarily}
choose the expensive resort over the cheap, crowded one, and the poor must not invade the
premium tier. The same distaste $\aHL$ that creates surplus also relaxes the high type's
incentive constraint. The asymmetry is the screening device.
\item \textbf{Rent accounting.} A single posted price must leave a rent to one type (the
pooling rent $\TP$). Sorting recovers $\TP$ but must pay whatever screening rent $R$ is needed
to separate the types. The profit comparison nets these against the concentration cost.
\end{enumerate}

\paragraph{Relation to the literature.}
The model sits at the intersection of three strands. Congestion and population games
(Rosenthal, 1973; Sandholm, 2010) and player--specific congestion (Milchtaich, 1996) supply
the equilibrium notion but reduce the externality to scalar load. Club--good theory (Buchanan,
1965; Scotchmer, 1994, 2002) captures composition preferences but is static and non--strategic
on the provider side. Schelling (1971) has composition thresholds but no prices and no firms.
Our contribution is to combine \emph{composition--dependent} payoffs with an \emph{asymmetric}
cross--type derivative \emph{and} a strategic (here, colluding) provider layer, and to
characterize exactly when type--segregation is profitable and stable.

% ======================================================================
\section{Primitives}
% ======================================================================

\paragraph{Population.}
There is a nonatomic continuum of tourists of two types $\theta \in \{H,L\}$, with total
masses $\mH>0$ (rich) and $\mL>0$ (poor); write $M=\mH+\mL$. A type--$\theta$ tourist has
gross value (budget, willingness to pay) $v_\theta>0$, with a strict value gap
\begin{equation}
\vH > \vL > 0, \qquad d \coloneqq \vH-\vL .
\end{equation}
Each tourist chooses one resort or an outside option normalized to utility $0$.

\paragraph{Resorts.}
There are $J\ge 2$ resorts, indexed $j\in\{1,\dots,J\}$, \emph{intrinsically identical}.
Resort $j$ posts a price $p_j\ge 0$; marginal cost is normalized to $0$. Let $x_{\theta j}\ge0$
be the mass of type $\theta$ at resort $j$ and $x_{\theta 0}\ge 0$ the mass taking the outside
option, with adding--up $x_{\theta 0}+\sum_j x_{\theta j}=m_\theta$.

\paragraph{Payoffs.}
A type--$\theta$ tourist at resort $j$ obtains
\begin{equation}
u_{\theta j}(x,p) \;=\; v_\theta - p_j - c_\theta\!\big(x_{Hj},x_{Lj}\big),
\label{eq:payoff}
\end{equation}
where the congestion / peer cost is linear in the composition,
\begin{equation}
\begin{pmatrix} c_H \\ c_L \end{pmatrix}
=
\begin{pmatrix} \aHH & \aHL \\ \aLH & \aLL \end{pmatrix}
\begin{pmatrix} x_{Hj} \\ x_{Lj} \end{pmatrix},
\qquad
A=\begin{pmatrix}\aHH & \aHL\\ \aLH & \aLL\end{pmatrix}.
\label{eq:congestion}
\end{equation}
The cartel maximizes joint profit $\Pi(p,x)=\sum_j p_j\,(x_{Hj}+x_{Lj})$ and may use transfers
across resorts. Because resorts are identical, any asymmetric policy must be sustained by
prices and the resulting composition, not by an exogenous quality difference.

\begin{assumption}[Asymmetric peer externality]\label{ass:asym}
The rich dislike crowding by the poor far more than the reverse, $\aHL \gg \aLH \ge 0$. The
clean benchmark is the one--directional case $\aLH=0$. Own--type crowding is strictly positive,
$\aHH,\aLL>0$, capturing physical congestion (lift lines, restaurants) and preventing
degenerate concentration. The aspirational variant $\aLH<0$ is treated in
Section~\ref{sec:discussion}.
\end{assumption}

% ======================================================================
\section{Tourist equilibrium given prices}
% ======================================================================

Fix a price profile $\bm{p}=(p_1,\dots,p_J)$. Tourists are followers; resorts are leaders (a
Stackelberg structure). Follower play is a population game.

\begin{definition}[Wardrop--Nash equilibrium]\label{def:eq}
A feasible allocation $x$ is an equilibrium given $\bm p$ if for each type $\theta$ there is a
scalar $\bar u_\theta\ge 0$ such that
\[
x_{\theta j}>0 \Rightarrow u_{\theta j}(x,p)=\bar u_\theta,
\qquad
x_{\theta j}=0 \Rightarrow u_{\theta j}(x,p)\le \bar u_\theta,
\qquad
x_{\theta 0}>0 \Rightarrow \bar u_\theta=0.
\]
Each type uses only utility--maximizing resorts and takes the outside option only when no
resort yields positive utility.
\end{definition}

\begin{remark}[Not a potential game in general]\label{rem:nopotential}
Since tourists are nonatomic, a deviator treats resort loads as fixed; with continuous costs an
equilibrium exists. The follower game admits an exact potential only when the cross--effects are
symmetric, $\aHL=\aLH$; under Assumption~\ref{ass:asym} the matrix $A$ is asymmetric, so the
game is \emph{not} a potential game and equilibria may be non--unique. We therefore work with
explicit allocations and verify equilibrium by the inequalities of Definition~\ref{def:eq}
rather than by appeal to a potential. When all resorts in a tier share price and composition
and $\aHH,\aLL>0$, own--type congestion equalizes loads within the tier.
\end{remark}

% ======================================================================
\section{Baseline: symmetric pricing (pooling)}\label{sec:pool}
% ======================================================================

Suppose every resort posts the same price $p$. The symmetric pooled allocation
\begin{equation}
x_{Hj}^{P}=\frac{\mH}{J},\qquad x_{Lj}^{P}=\frac{\mL}{J}\qquad\text{for all }j
\end{equation}
is an equilibrium whenever both types participate. Define each type's net (congestion--adjusted)
willingness to pay:
\begin{equation}
\wHP = \vH - \frac{\aHH\mH+\aHL\mL}{J},
\qquad
\wLP = \vL - \frac{\aLH\mH+\aLL\mL}{J}.
\label{eq:wpool}
\end{equation}
A single posted price cannot separate types. Assuming $\min\{\wHP,\wLP\}\ge0$, the largest
common price keeping both types in the market is $p^{P}=\min\{\wHP,\wLP\}$, so
\begin{equation}
\boxed{\ \PiP = M\,p^{P} = M\min\{\wHP,\wLP\}. \ }
\label{eq:Ppool}
\end{equation}
Total tourist rent under pooling and total surplus are
\begin{align}
\TP &= \mH(\wHP-p^{P})+\mL(\wLP-p^{P}) \;\ge 0, \label{eq:Tpool}\\
\WP &= \mH\vH+\mL\vL-\frac{\aHH\mH^2+(\aHL+\aLH)\mH\mL+\aLL\mL^2}{J},
\end{align}
and since marginal cost is zero, $\PiP=\WP-\TP$. A common price leaves a positive rent to the
type whose net WTP is not the minimum; this recovered rent reappears below. (Equation
\eqref{eq:Ppool} is the pooling benchmark; a symmetric cartel could instead \emph{exclude} one
type, e.g.\ serve only $H$ at profit $\mH(\vH-\aHH\mH/J)$ if that price also deters $L$. A claim
that sorting beats \emph{every} symmetric policy must also dominate these exclusion branches; we
return to this in Remark~\ref{rem:exclusion}.)

% ======================================================================
\section{Coordinated differentiation (separating regime)}\label{sec:sort}
% ======================================================================

The cartel designates $K\in\{1,\dots,J-1\}$ resorts as \emph{premium} at price $p_H$ and
$J-K$ resorts as \emph{budget} at price $p_L$. The intended separating allocation is all rich
at premium resorts and all poor at budget resorts, evenly split within each tier:
\begin{equation}
\text{premium: } x_H=\frac{\mH}{K},\ x_L=0;
\qquad
\text{budget: } x_H=0,\ x_L=\frac{\mL}{J-K}.
\end{equation}
Define the four congestion magnitudes
\begin{equation}
h_K=\frac{\aHH\mH}{K},\quad
\ell_K=\frac{\aLL\mL}{J-K},\quad
z^{HL}_K=\frac{\aHL\mL}{J-K},\quad
z^{LH}_K=\frac{\aLH\mH}{K},
\end{equation}
where $h_K,\ell_K$ are on--path own--type costs, $z^{HL}_K$ is the cost a rich deviator would
incur at a budget resort, and $z^{LH}_K$ the cost a poor deviator would incur at a premium
resort.

\subsection{Incentive and participation constraints}

A deviation is infinitesimal, so it does not change a resort's composition. The separating
allocation is an equilibrium iff
\begin{alignat}{2}
&\text{(IC--H)} &\qquad p_H+h_K &\le p_L+z^{HL}_K, \label{eq:ICH}\\
&\text{(IC--L)} &\qquad p_L+\ell_K &\le p_H+z^{LH}_K, \label{eq:ICL}\\
&\text{(IR--H)} &\qquad p_H+h_K &\le \vH, \label{eq:IRH}\\
&\text{(IR--L)} &\qquad p_L+\ell_K &\le \vL. \label{eq:IRL}
\end{alignat}
Define the \emph{generalized prices} (price plus own congestion borne) and two key gaps,
\begin{equation}
q_H=p_H+h_K,\quad q_L=p_L+\ell_K,
\qquad
A_K=z^{HL}_K-\ell_K=\frac{(\aHL-\aLL)\mL}{J-K},
\qquad
B_K=h_K-z^{LH}_K=\frac{(\aHH-\aLH)\mH}{K}.
\end{equation}
The two incentive constraints collapse to the clean sandwich
\begin{equation}
\boxed{\ B_K \ \le\ q_H-q_L \ \le\ A_K. \ }
\label{eq:sandwich}
\end{equation}
$A_K$ is the largest generalized premium the rich will accept over the budget tier; a larger
$\aHL$ raises $A_K$, relaxing \eqref{eq:ICH}. $B_K$ is the smallest premium that keeps the poor
out. Price--only separation for a given $K$ is feasible iff
\begin{equation}
A_K\ge B_K
\qquad\Longleftrightarrow\qquad
z^{HL}_K+z^{LH}_K\ \ge\ h_K+\ell_K .
\label{eq:feasible}
\end{equation}

\begin{remark}[The overlooked constraint]\label{rem:ICL}
Strong one--way aversion (large $\aHL$) helps deter the \emph{rich} from the budget tier
through $A_K$, but it does nothing for $B_K$. Worse, the asymmetry works against
\eqref{eq:ICL}: a \emph{small} $\aLH$ \emph{raises} $B_K=(\aHH-\aLH)\mH/K$, so the poor find
the premium tier relatively attractive and must be priced out. One--way aversion is therefore
helpful through $\aHL$ but not mechanically helpful for every constraint; feasibility requires
checking \eqref{eq:feasible}, not merely a large $\aHL$.
\end{remark}

\subsection{Optimal screening prices for a fixed split}

\begin{proposition}[Optimal separating prices]\label{prop:prices}
Fix an \emph{admissible} $K$ (one with $A_K\ge B_K$ and nonnegative prices below). The
joint--profit--maximizing generalized prices are
\begin{equation}
(q_H^\star,q_L^\star)=
\begin{cases}
(\vL+A_K,\ \vL), & d>A_K,\\[2pt]
(\vH,\ \vL), & B_K\le d\le A_K,\\[2pt]
(\vH,\ \vH-B_K), & d<B_K,
\end{cases}
\qquad
p_H^\star=q_H^\star-h_K,\quad p_L^\star=q_L^\star-\ell_K,
\end{equation}
and the resulting total tourist (screening) rent is
\begin{equation}
\boxed{\ R_K=\mH\,\pos{d-A_K}+\mL\,\pos{B_K-d}, \qquad \pos{z}=\max\{z,0\}. \ }
\label{eq:rent}
\end{equation}
\end{proposition}

\begin{proof}
For fixed $K$ the costs and quantities are constants, so the cartel maximizes
$\mH(q_H-h_K)+\mL(q_L-\ell_K)$ subject to $q_H\le\vH$, $q_L\le\vL$ and \eqref{eq:sandwich}. The
objective increases in both generalized prices, so push each up to its binding constraint. If
$d>A_K$ the two IR constraints cannot both bind (their gap $d$ exceeds the allowed $A_K$), so
$q_L=\vL$ and \eqref{eq:ICH} binds at $q_H=\vL+A_K$. If $B_K\le d\le A_K$ both IR bind,
$(q_H,q_L)=(\vH,\vL)$. If $d<B_K$ then $q_H=\vH$ and \eqref{eq:ICL} binds at $q_L=\vH-B_K$.
Tourist rent is $\sum_\theta m_\theta(v_\theta-q_\theta^\star)$, which gives \eqref{eq:rent}.
\end{proof}

\begin{lemma}[The asymmetry is the screening device]\label{lem:screen}
In the middle (full--extraction) case $B_K\le d\le A_K$ the differentiated prices implement
exactly what type--contingent pricing would: resort choice reveals type and the cartel extracts
all surplus, $R_K=0$. The high type's information rent is $\mH\pos{d-A_K}$, which falls one for
one as $\aHL$ rises (through $A_K$) until it vanishes. This is the precise sense in which
asymmetric peer congestion is a screening instrument.
\end{lemma}

\subsection{Sorted surplus and profit}
With cross--type contact eliminated, total surplus and profit are
\begin{equation}
\WS(K)=\mH\vH+\mL\vL-\frac{\aHH\mH^2}{K}-\frac{\aLL\mL^2}{J-K},
\qquad
\boxed{\ \PiS(K)=\WS(K)-R_K. \ }
\label{eq:Psort}
\end{equation}

% ======================================================================
\section{Main result: when does sorting beat pooling?}\label{sec:main}
% ======================================================================

The real--surplus gain from sorting is
\begin{equation}
\Delta W(K)=\WS(K)-\WP
=\underbrace{\frac{(\aHL+\aLH)\mH\mL}{J}}_{\text{cross--type relief}}
-\underbrace{\Big[\aHH\mH^2\big(\tfrac1K-\tfrac1J\big)+\aLL\mL^2\big(\tfrac1{J-K}-\tfrac1J\big)\Big]}_{\text{own--type concentration cost}} .
\label{eq:dW}
\end{equation}

\begin{proposition}[Exact sorting--dominance condition]\label{prop:main}
For any admissible $K$, coordinated sorting strictly raises joint profit over the symmetric
pooling benchmark, $\PiS(K)>\PiP$, if and only if
\begin{equation}
\boxed{\ \Delta W(K) \;+\; \TP \;-\; R_K \;>\;0. \ }
\label{eq:maincond}
\end{equation}
\end{proposition}

\begin{proof}
$\PiP=\WP-\TP$ from \eqref{eq:Ppool}--\eqref{eq:Tpool} and $\PiS(K)=\WS(K)-R_K$ from
\eqref{eq:Psort}. Subtract: $\PiS(K)-\PiP=\big(\WS(K)-\WP\big)+\TP-R_K=\Delta W(K)+\TP-R_K$.
\end{proof}

\noindent Equation \eqref{eq:maincond} is transparent: sorting wins when the congestion relief
plus the recovered pooling rent $\TP$ exceeds the own--type concentration cost plus the
screening rent $R_K$ it must pay. The asymmetric term $\aHL$ enters \emph{twice}: directly in
the relief term of $\Delta W$, and indirectly by shrinking $R_K$ through $A_K$
(Lemma~\ref{lem:screen}). It is \emph{not} enough that $\aHL$ be large: a severe own--congestion
($\aHH,\aLL$ large) or an infeasible split ($A_K<B_K$, Remark~\ref{rem:ICL}) can still defeat
sorting.

\subsection{Optimal number of premium resorts}

Within the full--extraction region ($R_K=0$), only the concentration cost depends on $K$, so the
cartel minimizes $C(K)=\aHH\mH^2/K+\aLL\mL^2/(J-K)$.

\begin{corollary}[Optimal split and clean threshold]\label{cor:split}
The interior minimizer is
\begin{equation}
\frac{K^\star}{J}=\frac{\sqrt{\aHH}\,\mH}{\sqrt{\aHH}\,\mH+\sqrt{\aLL}\,\mL},
\qquad
C(K^\star)=\frac{\big(\sqrt{\aHH}\,\mH+\sqrt{\aLL}\,\mL\big)^2}{J},
\end{equation}
rounded to the nearest admissible integer. If $K^\star$ lies in the full--extraction region
($B_{K^\star}\le d\le A_{K^\star}$), then
\begin{equation}
\boxed{\ \PiS(K^\star)-\PiP=\frac{\big(\aHL+\aLH-2\sqrt{\aHH\,\aLL}\big)\,\mH\mL}{J}+\TP, \ }
\end{equation}
so sorting strictly dominates pooling iff that expression is positive. Comparing total
\emph{surplus} (ignoring the rent transfer $\TP$) gives the especially clean criterion
\begin{equation}
\boxed{\ \aHL+\aLH \;>\; 2\sqrt{\aHH\,\aLL}. \ }
\label{eq:surpluscond}
\end{equation}
\end{corollary}

\noindent Under near one--way aversion ($\aLH\approx0$), a sufficiently large $\aHL$ alone
satisfies \eqref{eq:surpluscond}. The extra $\TP\ge0$ shows the cartel can profit from sorting
even when sorting slightly \emph{lowers} total surplus, because it also improves price
discrimination.

\begin{remark}[Dominating exclusion policies]\label{rem:exclusion}
To claim sorting beats \emph{every} symmetric policy, not just pooling, one adds
$\PiS(K)>\mH(\vH-\aHH\mH/J)$ and the analogous $L$--only inequality. These are easy to check
numerically and hold whenever the served mass and value gap are large enough that excluding a
whole type is wasteful.
\end{remark}

% ======================================================================
\section{The role of the asymmetry}\label{sec:role}
% ======================================================================

The asymmetry has two distinct roles and one trap.
\begin{enumerate}[label=(\alph*),leftmargin=2.2em]
\item \textbf{Surplus source.} $\aHL$ enters the relief term $(\aHL+\aLH)\mH\mL/J$ in
\eqref{eq:dW}. With $\aLH\approx0$, separating the rich from the poor removes a real loss for
the rich without removing a corresponding benefit for the poor.
\item \textbf{Screening technology.} $\aHL$ raises $A_K=(\aHL-\aLL)\mL/(J-K)$, relaxing the high
type's incentive constraint \eqref{eq:ICH} and shrinking the rent $R_K$ in \eqref{eq:rent}. The
budget tier, crowded with poor, is an unattractive outside option for the rich.
\item \textbf{The countervailing trap.} A \emph{low} $\aLH$ \emph{raises}
$B_K=(\aHH-\aLH)\mH/K$, the premium the cartel must charge to keep the poor out
(Remark~\ref{rem:ICL}). If $B_K$ exceeds $A_K$ the separating menu is infeasible at any prices.
So the very feature that makes the rich willing to pay can make the poor eager to follow, and
both incentive bounds must be checked.
\end{enumerate}

% ======================================================================
\section{What changes or breaks the result}\label{sec:discussion}
% ======================================================================

\paragraph{Symmetric externality ($\aHL=\aLH>0$).}
Cross--type contact is mutually costly; sorting yields a larger surplus gain and both types
weakly prefer separation. The mechanism is no longer elite avoidance but mutual segregation.
The fully homogeneous case $\aHH=\aHL=\aLH=\aLL$ gives $A_K=B_K=0$ at the equal split: resort
labels carry no screening content and any apparent gain vanishes once the rent is accounted.

\paragraph{Aspirational poor ($\aLH<0$).}
If the poor enjoy proximity to the rich, sorting destroys that benefit: the relief term
$(\aHL+\aLH)\mH\mL/J$ shrinks and $B_K=(\aHH-\aLH)\mH/K$ rises, hurting sorting through both the
surplus and the IC channels. For sufficiently negative $\aLH$, price--only separation can be
infeasible or unprofitable even when the rich strongly dislike the poor. This is the empirically
interesting ``see and be seen'' case.

\paragraph{Large own--type congestion.}
Sorting concentrates each type into fewer resorts. When $\aHH$ or $\aLL$ is large the
concentration cost in \eqref{eq:dW} can exceed the cross--type relief; this is the central
technological trade--off and is exactly criterion \eqref{eq:surpluscond}.

\paragraph{Failure of price--only screening.}
If $A_K<B_K$ for every $K$, no pair of posted prices supports full separation: the price needed
to deter the poor pushes the rich to the budget tier. Partial mixing, exclusion, or a second
instrument (capacity, branding, minimum spend) is then required.

\paragraph{Cartel enforceability.}
The analysis maximizes joint profit with transfers. Without transfers, premium and budget
resorts earn different profits and may disagree on the policy. Without enforcement, a single
resort may deviate to attract both types. The non--cooperative pricing game among resorts is a
separate question; notably the asymmetric congestion is self--limiting there, since a premium
resort that cuts price to add volume imports the poor and the very cross--congestion that
repels its rich clientele.

% ======================================================================
\section{Dynamics and the wrong--equilibrium trap}\label{sec:dynamics}
% ======================================================================

The pooled allocation is an equilibrium whenever all resorts post one price and both types
participate, \emph{even when} a sorted policy yields strictly higher profit and surplus. Since
resorts are identical, moving from pooling to sorting requires a coordinated change in prices or
expectations. Consider the within--type replicator dynamic
$\dot y_{\theta j}=y_{\theta j}\,(u_{\theta j}-\bar u_\theta)$ with $y_{\theta j}=x_{\theta j}/m_\theta$;
Wardrop equilibria are its rest points. Perturbing two equal--price resorts by moving masses
$(\delta_H,\delta_L)$ from one to the other, the payoff gap for type $\theta$ is
$u_{\theta 1}-u_{\theta 2}=-2(\alpha_{\theta H}\delta_H+\alpha_{\theta L}\delta_L)$, so the
linearization is governed by $-\operatorname{diag}(\mH,\mL)\,A$. The pooled state is locally
stable iff
\begin{equation}
\boxed{\ \aHH\,\aLL \;>\; \aHL\,\aLH. \ }
\label{eq:stability}
\end{equation}
Under strong one--way aversion ($\aLH\approx0$) the right side is near zero, so pooling stays
locally stable \emph{however large $\aHL$ is}, even though sorting is more profitable and, by
\eqref{eq:surpluscond}, more efficient. This is a sharp wrong--equilibrium trap: unilateral
tourist adjustment does not create the sorted outcome, and only a coordinated policy jump (price
plus, if needed, capacity) crosses the basin boundary. A sorted state with \emph{strict}
cross--tier incentives is locally resistant to invasion; when an IC bound binds it is only
weakly resistant and small shocks can reintroduce mixing.

% ======================================================================
\section{Extensions}\label{sec:ext}
% ======================================================================

\paragraph{Capacity.}
A cap $k_j$ changes equilibrium congestion and can ration the premium tier, repairing the
invasion failure $A_K<B_K$; but a cap alone does not distinguish types, so it screens only when
paired with a price (or auction, minimum spend, priority booking) that selects high WTP.

\paragraph{Branding and amenities.}
If resort $j$ chooses a brand $b_j$ shifting gross values to $v_\theta(b_j)$, a premium brand
supports sorting by raising the value gap $d=\vH(b_j)-\vL(b_j)$, moving the economy into the
full--extraction window $B_K\le d\le A_K$.

\paragraph{Hard budgets.}
If the poor face a hard budget $\bar p_L$, then $p_H>\bar p_L$ mechanically blocks invasion,
relaxing or removing \eqref{eq:ICL}. The present model is more demanding: separation works
through preferences, not an exogenous inability to pay.

\paragraph{More than two types.}
With a continuum of types resorts form a menu of (price, peer--environment) pairs; the two--type
sandwich \eqref{eq:sandwich} becomes adjacent--type incentive constraints, and a single--crossing
condition on type--specific peer costs delivers monotone sorting.

% ======================================================================
\section{Conclusion}
% ======================================================================

Identical resorts, a heterogeneous population, and an asymmetric peer externality
($\aHL\gg\aLH$) suffice to make coordinated price differentiation strictly more profitable than
symmetric pricing, under the exact condition $\Delta W(K)+\TP-R_K>0$ of
Proposition~\ref{prop:main}. Segregation removes the congestion the rich suffer from the poor;
because the poor are nearly indifferent the relief is largely pure surplus, which the cartel
captures by charging the rich a premium. The asymmetry's deeper role is as a screening
technology: $\aHL$ relaxes the high type's incentive constraint and shrinks its information
rent. Yet asymmetry is not a free lunch. A low $\aLH$ makes the poor harder to exclude, the
concentration cost can dominate, and (criterion \eqref{eq:stability}) the profitable sorted
outcome can be dynamically unreachable from pooling without coordinated policy, the
wrong--equilibrium trap. At the optimal split the efficiency of sorting reduces to the single
inequality $\aHL+\aLH>2\sqrt{\aHH\aLL}$.

\appendix
\section{Notation}
\begin{center}
\begin{tabular}{ll}
\toprule
Symbol & Meaning \\
\midrule
$H,L$ & rich/high--value and poor/low--value tourist types \\
$\mH,\mL,\,M$ & mass of each type; total $M=\mH+\mL$ \\
$\vH,\vL,\,d$ & gross values, $\vH>\vL$; value gap $d=\vH-\vL$ \\
$J$ & number of identical resorts \\
$K,\,J-K$ & number of premium and budget resorts \\
$p_j$ & price posted by resort $j$ \\
$x_{\theta j}$ & mass of type $\theta$ at resort $j$ \\
$\aHH,\aLL$ & own--type (physical) congestion sensitivities \\
$\aHL$ & rich aversion to poor crowding (large) \\
$\aLH$ & poor aversion to rich crowding ($\approx0$, or $<0$ aspirational) \\
$h_K,\ell_K$ & on--path own--type congestion costs \\
$z^{HL}_K,z^{LH}_K$ & cross--tier deviation congestion costs \\
$q_H,q_L$ & generalized prices, $q_\theta=p_\theta+(\text{own congestion})$ \\
$A_K,B_K$ & upper / lower bound on the generalized premium $q_H-q_L$ \\
$R_K$ & screening rent left to tourists under sorting \\
$\TP$ & rent a single pooling price must leave \\
$W^{\bullet},\Pi^{\bullet}$ & total surplus and joint (cartel) profit \\
\bottomrule
\end{tabular}
\end{center}

\end{document}
